\section{Background and Related Work}
\label{sec:background}

We synthesize four strands of literature that bracket the contribution
--- LOB deep learning, DeFi lending mechanics, ML-based DeFi rate
forecasting, and MCDM in algorithmic allocation --- and then position
the work against the recent impossibility result of
\cite{halpern2024fair}, the most relevant theoretical counter-evidence.

\subsection{LOB microstructure deep learning}
\label{sec:bg-lob}

Modern LOB mid-price prediction has converged on convolutional-recurrent
backbones. DeepLOB \cite{zhang2019deeplob} established the
CNN+LSTM template on the FI-2010 benchmark \cite{ntakaris2018benchmark};
\cite{sirignano2019universal} demonstrated cross-stock universal features
under deep learning; \cite{tsantekidis2017using, briola2021deep,
tran2018temporal} extended the family with attention and bilinear
mechanisms. The author's prior dual-branch architecture
\cite{solovev2026lob} is the immediate predecessor to the present work
and we recap its key results in \Cref{sec:lob-recap}. Recent crypto-LOB
work \cite{briola2025microstructure, lit2025limit} reinforces a finding
central to our motivation: \emph{domain-aware input engineering
dominates raw depth scaling}, supporting the
two-branch-over-deeper-stack design.

\subsection{DeFi lending mechanics}
\label{sec:bg-defi-mechanics}

Aave V3 and Compound V3 USDC supply rates are deterministic
piecewise-linear functions $f_{\text{kink}}(u)$ of pool utilization
$u = \text{borrowed} / \text{supplied}$, with a shallow slope below a
kink point (typically $u^{*} = 0.8$--$0.9$) and a steep slope above it
that taxes high utilization. The strategy parameters are protocol
governance variables: Aave's \texttt{ReserveInterestRateStrategy}
contract and Compound V3's hard-coded slopes. Although the function
$f_{\text{kink}}$ is deterministic given $u$, the realised rate stream
is stochastic because $u$ moves with on-chain demand.

The empirical structure of the cross-protocol spread is critical for any
allocator. \cite{gudgeon2020loanable} establishes cointegration between
Aave and Compound USDC rates with Compound leading and Aave adjusting at
speed $0.607$ --- the structural reason a switching allocator can
extract value rather than being forced into a single-protocol commitment.
Recent crosschain comparative analyses \cite{aavecompoundcrosschain2025}
confirm the cointegration persists into 2025.  Our 12,895-hour panel
(\Cref{sec:defi-experiment}) finds the spread structure changes
regime around 2025-Q3 $\to$ Q4, with the Aave-pays-more share rising from
$39.9\%$ to $56.3\%$ across that quarter boundary.

\subsubsection{Market structure (April 2026 snapshot)}
\label{sec:bg-market-structure}

The DeFi lending market is highly concentrated.  As of April 2026, total
deposits across all lending protocols total approximately
\$54\,B, with the top ten protocols accounting for roughly $78\%$ of
that value \cite{defillama2026snapshot}.  \Cref{tab:lending-tvl} reports
the top-eight breakdown; the two protocols highlighted in bold delimit
the scope of our empirical study.

\begin{table}[t]
  \centering
  \caption{Top-eight DeFi lending protocols by total value locked (TVL),
           April 2026 \cite{defillama2026snapshot}.  Bold rows mark the
           subset covered by this paper's empirical study.}
  \label{tab:lending-tvl}
  \begin{tabular}{lrr}
    \toprule
    Protocol (chain)            & TVL (\$B) & Share \\
    \midrule
    \textbf{Aave V3}            & \textbf{19.4} & \textbf{36.0\%} \\
    Spark (SparkLend)           & 6.8           & 12.6\% \\
    Morpho Blue                 & 4.9           & 9.1\%  \\
    \textbf{Compound V3}        & \textbf{2.7}  & \textbf{5.0\%} \\
    JustLend (Tron)             & 2.4           & 4.4\%  \\
    Fluid                       & 1.6           & 3.0\%  \\
    Kamino (Solana)             & 1.1           & 2.0\%  \\
    Euler V2                    & 0.89          & 1.6\%  \\
    \midrule
    Top-8 subtotal              & 39.8          & 73.7\% \\
    Tail ($\sim 350+$ protocols)& 14.2          & 26.3\% \\
    \bottomrule
  \end{tabular}
\end{table}

Our empirical study covers Aave V3 and Compound V3, which jointly hold
\$22.1\,B (\$19.4\,B + \$2.7\,B), or $41\%$ of total lending TVL.  This
is a representative, dominant subset of the market rather than a niche.
The DA-BiGRU-CNN forecaster and the 4-factor MCDM allocator generalize
to $N$-way protocol allocation without architectural change: each
additional protocol requires only its protocol-specific kink function
$f_{\text{kink}}^{(i)}(u)$ and a per-protocol Branch~B feature wiring.
The methodology therefore applies unchanged to the remaining top-eight
protocols in \Cref{tab:lending-tvl}; we leave the $N$-way evaluation as
future work (\Cref{sec:limitations}).

\subsection{Forecasting in DeFi rates}
\label{sec:bg-forecasting}

Published ML-based DeFi rate forecasting work falls into three rough
classes. (i) \emph{Protocol-side rate control.}
\cite{bertucci2025rl} train offline reinforcement-learning agents (CQL,
BC, TD3-BC) on historical Aave data to set rates from the protocol
governance side; the closely related Auto.gov framework
\cite{xu2023autogov} reports $\sim14\%$ improvement over benchmarks for
parameter governance under RL. (ii) \emph{Adaptive rate control.}
\cite{agilerate2024} fits recursive least-squares (RLS) controllers on
Compound demand/supply curves at 3-hour resolution. (iii) \emph{Optimal
protocol design.} \cite{baude2025optimal} derives a closed-form
Riccati-ODE solution for the protocol-optimal rate curve under linear
agent response and a Monte-Carlo + deep-learning estimator for the
nonlinear case; block-level empirical calibration shows industry-standard
piecewise rate curves are sub-optimal, leaving structural residuals.
This last point is foundational for us: \cite{baude2025optimal} solves
the \emph{protocol designer's} problem (which rate function should the
contract use?), whereas we solve the \emph{lender's} problem (given the
realised rate stream, where should liquidity go?). Our DA-BiGRU-CNN
forecaster targets exactly the structural residuals their optimality
analysis predicts.  We are not aware of a published forecast-driven
\emph{multi-protocol MCDM allocator} in the May 2026 literature; the
nearest neighbour is the LSTM+Q-learning AMM quoter of
\cite{predamm2024}, which transplants a TradFi forecaster to Uniswap V3
but does not span multiple lending protocols nor use an MCDM head.

\subsection{MCDM in algorithmic allocation}
\label{sec:bg-mcdm}

Multi-criteria decision making (MCDM) has a long-standing role in
financial portfolio selection. TOPSIS \cite{hwang1981mcdm} and PROMETHEE
\cite{brans1985promethee} are the two dominant families. Recent
cryptocurrency-specific applications include
\cite{aljinovic2021crypto}'s multicriteria crypto portfolio selection
and \cite{hosseinzadeh2023promethee}'s asymmetric PROMETHEE-II that
integrates return prediction with multi-criteria ranking. Our 4-factor
TOPSIS-style allocator (APY, Risk, Cost, Stability) is in this tradition
and inherits its criteria-and-weights structure from the author's
ERC-4626 yield-vault project \cite{solovev2026vault}, modifying it by
substituting forecast-driven APY for EMA-smoothed spot APY.

\subsection{Counter-evidence: Halpern--Pass--Saraf impossibility}
\label{sec:bg-impossibility}

The most pointed theoretical counter-evidence to forecast-driven DeFi
lending strategies is \cite{halpern2024fair}.  Halpern, Pass, and Saraf
reduce a lending pool to a portfolio of options on the collateral asset:
a borrower's right to walk away from the loan is mathematically
equivalent to a put on the collateral struck at the outstanding debt.
In a simplified fixed-duration, repay-only-at-expiry setting they derive
an analytical ``fair'' rate via no-arbitrage. They then prove that in
the realistic setting --- dynamic, perpetual loans with early repayment
of the kind Aave V3 and Compound V3 implement --- \emph{no fair
interest rate exists}: any rate either over-compensates lenders relative
to the option-implied premium or under-compensates them, because the
borrower's repayment optionality has no finite hedging cost.  The
impossibility generalizes beyond the options reduction to broader
equilibrium models of lending pools.

A naive reading of \cite{halpern2024fair} suggests there is no
exploitable structure in DeFi rates and any allocator should
underperform a passive benchmark. We argue the implication does not
follow.  Halpern et al.\ prove a \emph{static equilibrium fairness}
impossibility under no-arbitrage --- a statement about whether a
single rate function can correctly price the borrower-side optionality
in equilibrium.  Our claim is the orthogonal \emph{empirical
predictability} statement: the realised rate stream observed off-chain
exhibits enough short-horizon structure (cointegration, regime
persistence, kink-residual cycles) to drive a forecast-and-allocate
strategy that improves a multi-criteria score against a reactive
baseline.  Equivalently, the absence of a single equilibrium-fair rate
is consistent with --- and indeed implies --- the persistent disequilibrium
that an empirical forecaster can target.  The two claims constrain
different layers of the protocol stack: \cite{halpern2024fair} constrains
pricing theory at the contract design layer; we operate at the
lender-side execution layer.  Reporting the honest H0-publishable result
(\Cref{sec:limitations}) makes this distinction empirically verifiable
rather than rhetorical.
